\section{Exact compilers and the limits of representation-specific rank}
\label{sec:exact-compilers}

A formulation can become smaller because several input constraints have
one exact replacement, because a source-dependent function has a finite
summary, or because auxiliary variables compose simpler functions. These
are different operations. We give precise criteria for the first two and
show why exponential-feature rank does not obstruct the third. Throughout
this section, a \emph{compiler} means a map from the specified input data
to an exact representation with the specified output format. A restriction
on that format is part of the theorem.

\subsection{Intersections of translates}

For a proper cone $K\subset\R^d$, write $x\leq_K y$ if $y-x\in K$.
The following is the finite-dimensional vector-lattice characterization
of simplicial cones, expressed as an exact compilation criterion. The
order-theoretic result is classical; see, for example,
\citet[Introduction]{Nemeth2004} and the references there.

\begin{proposition}[One-translate compilation]\label{prop:cone-meet}
Let $d\geq1$. The identity
\[
 (a-K)\cap(b-K)=c-K
\]
has a solution $c$ for every pair $a,b\in\R^d$ if and only if
$K=T\R_+^d$ for an invertible linear map $T$. In that case, for any finite
family,
\begin{equation}\label{eq:simplicial-compiler}
 \bigcap_{i=1}^N(a_i-K)=Tm-K,
 \qquad m_j=\min_i(T^{-1}a_i)_j.
\end{equation}
If the output apex must be one of the two input apices, the universal
pairwise property holds if and only if $d=1$.
\end{proposition}
\begin{proof}
The intersection identity says precisely that $c$ is the greatest lower
bound $a\wedge b$. Negation supplies least upper bounds, so the induced
order is a vector lattice. Here is a finite-dimensional proof of the
resulting simpliciality. Distinct extreme rays of $K$ are disjoint in the
lattice sense: a positive element below representatives of both rays
must lie on both and therefore vanish. The lattice decomposition identity
\[
 0\leq z\leq x+y
 \quad\Longrightarrow\quad
 z=(z\wedge x)+(z-x)^+,
 \quad 0\leq z\wedge x\leq x,
 \quad 0\leq(z-x)^+\leq y
\]
shows by induction that finite sums of disjoint groups of extreme rays
are disjoint. Any linear dependence among distinct extreme-ray
representatives would equate two nonzero positive sums over disjoint
groups, which is impossible. Thus there are at most $d$ extreme rays.
A compact base of a closed pointed cone is the convex hull of its extreme
points; since $K$ spans $\R^d$, there are exactly $d$ such rays and they
form a basis. This gives $K=T\R_+^d$. Coordinatewise minima prove the
converse and \eqref{eq:simplicial-compiler}.

An input apex can always be retained exactly when every pair is
comparable. This is equivalent to $K\cup(-K)=\R^d$. For $d\geq2$, the
unit sphere is connected, whereas its intersections with $K$ and $-K$
would be disjoint nonempty closed sets covering it. Pointedness gives
the contradiction. A one-dimensional proper cone has a total order.
\end{proof}

The number of incomparable input apices need not control the size of the
output: even an antichain of $N$ apices in a simplicial cone is summarized
by $d$ coordinate minima. For explicit classical output and coherent access
to an individual transformed entry $(T^{-1}a_i)_j$, quantum minimum
finding \cite{DurrHoyer1996} gives
\begin{equation}\label{eq:meet-queries}
 Q=O\bigl(d\sqrt N\log(d+1)\bigr),\qquad
 Q=\Omega(d\sqrt N),\qquad R=\Theta(dN).
\end{equation}
Here $Q$ and $R$ denote bounded-error quantum and randomized entry-query
complexities, respectively, and success means that all $d$ output
coordinates are correct. Exact comparison and retrieval of an encoded
entry are included in the oracle contract. For the upper bound, run each
minimum search with failure probability at most $1/(3d)$. For the quantum
lower bound, restrict each column to all ones or to a single hidden zero.
The output supplies $d$ independent search-decision answers, to which the
adversary direct-sum theorem applies \cite[Theorem~4]{ACGT2010}. For the
randomized lower bound, couple an input with independent equiprobable
all-one/single-zero columns to the input with any one column changed.
If an algorithm succeeds on the full answer vector with probability at
least $2/3$, each column must be distinguished with constant advantage;
a uniform hidden zero then forces $\Omega(N)$ expected probes in each
column. Summing over columns gives $\Omega(dN)$. The logarithm in the
quantum upper bound accounts for simultaneous success; no exact
$\Theta(d\sqrt N)$ upper bound is asserted here. If the supplied oracle
instead returns a physical coordinate of $a_i$ and $T^{-1}$ is dense,
forming one transformed entry can require $d$ source queries.

For the canonical logarithmic barriers, the scenario product has
parameter $Nd$ and the compiled translate has the optimal parameter $d$.
On identical-apex instances, a radial curve whose common transformed
slack is $\alpha\mathbf1$ has lengths
\[
 \sqrt{Nd}\left|\log(\alpha_1/\alpha_0)\right|
 \quad\hbox{and}\quad
 \sqrt d\left|\log(\alpha_1/\alpha_0)\right|,
\]
respectively. This follows directly from the Hessians of
$-N\sum_j\log s_j$ and $-\sum_j\log s_j$. It compares these barriers,
not all algorithms for the underlying optimization problem.

A factorwise radial restriction permits the same operation for arbitrary
proper cones. If $K=\prod_{j=1}^B K_j$, $e_j\in\intr K_j$, and
$a_i=(r_{i1}e_1,\ldots,r_{iB}e_B)$, then
\begin{equation}\label{eq:radial-product-compiler}
 \bigcap_i(a_i-K)=\prod_{j=1}^B(r_{*j}e_j-K_j),
 \qquad r_{*j}=\min_i r_{ij}.
\end{equation}
Indeed, the translates in each factor are nested in their scalar radius.
The corresponding query bounds are \eqref{eq:meet-queries} with $d=B$.
Given barriers of parameters $\nu_j$ on the factors, the compiled product
has a barrier of parameter $\sum_j\nu_j$, whereas the repeated scenario
barrier has parameter $N\sum_j\nu_j$. No symmetry of $K_j$ is needed.

\begin{proposition}[Lorentz-translate obstruction]
\label{prop:lorentz-translate-arcs}
For every $N$, there are $a_1,\ldots,a_N\in\R^3$ such that every finite
identity
\[
 \bigcap_{i=1}^N(a_i-Q_3)
   =\bigcap_{\ell=1}^m(c_\ell-Q_3)
\]
satisfies $m\geq N$.
\end{proposition}
\begin{proof}
Choose distinct $v_i\in\Sph^1$ and $0<\rho<R$, and set
$a_i=(0,\rho v_i)$. On the plane $t=-R$, the intersection is
$\bigcap_i B(\rho v_i,R)$. The point $z_i=-(R-\rho)v_i$ satisfies
\[
 \norm{z_i-\rho v_i}=R,
 \qquad
 \norm{z_i-\rho v_j}^2
 =(R-\rho)^2+\rho^2+2\rho(R-\rho)\ip{v_i}{v_j}<R^2
 \quad(j\ne i).
\]
Hence each original circle contributes an open boundary arc. An
alternative translate intersects this plane in a disk. Its radius must
be positive because the common intersection contains an open disk.
Every point of an original open arc lies on the boundary of at least
one alternative disk. Two distinct circles intersect in at most two
points, so finitely many alternative circles can cover that arc only
if one coincides with the original circle. Distinct original circles
therefore require distinct alternative disks.
\end{proof}

This obstruction already occurs in a cone of Jordan rank two. Copying
the three-dimensional variable at each constraint and imposing consensus
along a path gives a formulation whose graph of variable blocks is a
path. Thus a bound on cone rank or on this block graph does not bound
the surviving number of translates. The proposition concerns same-space,
same-oriented translates only. It gives no lower bound for lifts with
auxiliary variables, different cone maps, or custom barriers.

\subsection{Reusable summaries of perspective data}

Let $h:(0,\infty)\to\R$ be continuous. Its \emph{dilation rank} is
\[
 R_h=\dim\spanop\{r\mapsto h(\lambda r):\lambda>0\},
\]
possibly infinite. Consider the fixed-source perspective sum
\begin{equation}\label{eq:perspective-source}
 \Phi(U,V)=\sum_{i=1}^N w_i c_i U\,
       h\!\left(\frac{d_i}{c_i}\frac VU\right),
 \qquad U,V,w_i,c_i,d_i>0.
\end{equation}
A reusable summary must determine this function for every $U,V>0$;
it is not merely an answer at one query point.

\begin{theorem}[Continuous summary dimension]\label{thm:dilation-summary}
If $R_h=R<\infty$, the source dependence in
\eqref{eq:perspective-source} is summarized by $R$ scalar linear moments.
Conversely, choose any $q$ linearly independent dilates
$h(\lambda_1r),\ldots,h(\lambda_qr)$. If a continuous map
$E:\R_{++}^q\to\R^m$ admits a deterministic decoder recovering
$\sum_i a_i h(\lambda_i r)$ for every $r>0$ and every
$a\in\R_{++}^q$, then $m\geq q$. Consequently $R_h$ is the exact minimum
summary dimension for these unrestricted positive-coefficient families;
if $R_h=\infty$, no fixed finite dimension suffices.
\end{theorem}
\begin{proof}
Choose a basis $h_1,\ldots,h_R$ of the dilation space and write
$h(\lambda r)=\sum_\ell m_\ell(\lambda)h_\ell(r)$. Then
\[
 M_\ell=\sum_i w_i c_i m_\ell(d_i/c_i),
 \qquad
 \Phi(U,V)=U\sum_{\ell=1}^R M_\ell h_\ell(V/U).
\]
The coefficient functions $m_\ell$ are continuous: choose $R$ evaluation
points giving an invertible evaluation matrix and solve for the
coefficients using the continuous values $h(\lambda r_j)$.

For the lower bound, equality $E(a)=E(b)$ would give equality of the
reconstructed functions, hence $a=b$ by linear independence. Thus $E$
is injective. If $m<q$, append $q-m$ zeros to obtain a continuous
injection from the open set $\R_{++}^q$ into $\R^q$. Invariance of domain
\cite[Theorem~2B.3]{Hatcher2002} says its image is open, whereas the image
lies in a proper coordinate subspace. This is impossible.
\end{proof}

Decoder continuity is unnecessary. Encoder continuity and the
finite-dimensional real output are essential hypotheses; arbitrary
set-theoretic encodings into one infinite-precision real are excluded.
For a fixed restricted family of coefficients, the lower bound need
not equal $R_h$.

\begin{proposition}[Classical dilation classification]
\label{prop:dilation-classification}
The rank $R_h$ is finite if and only if
\begin{equation}\label{eq:dilation-exp-poly}
 h(r)=\sum_{j=1}^J r^{\rho_j}P_j(\log r),
\end{equation}
where the $\rho_j$ are distinct complex numbers and $P_j$ are nonzero
polynomials, with conjugate terms paired when $h$ is real-valued.
The empty sum represents $h=0$, with $R_h=0$. Otherwise
$R_h=\sum_j(\deg P_j+1)$.
\end{proposition}
\begin{proof}
Put $f(y)=h(e^y)$. Its translation span $W$ is finite dimensional
exactly when the dilation span is. Translation acts on $W$ by a
continuous group $T(t)$ of invertible matrices: continuity follows by
choosing an invertible finite evaluation matrix as above. For completeness,
choose small $\delta>0$ so that $M=\int_0^\delta T(t)\,dt$ is invertible.
The identity
$T(s)M=\int_s^{s+\delta}T(t)\,dt$ proves differentiability of $T$.
The group law then gives $T'(s)=T(s)A$ with $A=T'(0)$, so $T(s)=e^{sA}$.
Evaluation at zero shows that $f$ is an exponential polynomial; the
Jordan form of $A$ gives \eqref{eq:dilation-exp-poly}. Conversely,
translates of each $e^{\rho_j y}P_j(y)$ lie in the span of
$y^\ell e^{\rho_j y}$ for $0\leq\ell\leq\deg P_j$. These functions
are linearly independent. Each occurs in the cyclic translation span:
on this finite-dimensional span the differential operator acts with
one Jordan chain of length $\deg P_j+1$ for each distinct $\rho_j$,
and a polynomial in that operator isolates each chain. This also gives
the stated rank.
\end{proof}

The translation-space classification is due to Anselone--Korevaar
\cite{AnseloneKorevaar1964}; see also Engert's extension and historical
account \cite{Engert1970}. We included the elementary continuous
one-parameter proof to specify all assumptions. For example, $-\log r$
has rank two, $r^\alpha$ has rank one, and $e^r$ has infinite rank.
The latter follows from independence of distinct exponentials
$e^{\lambda r}$ on any open interval. The contribution of
Theorem~\ref{thm:dilation-summary} is the exact continuous-summary
interpretation, rather than a new classification of translation spaces.
It concerns source evaluation, not arbitrary conic lift size or barrier
parameter.

\subsection{A two-moment relative-entropy example}

For positive definite Hermitian matrices use the homogeneous, uncorrected
Umegaki convention
$D(X\|Y)=\tr X(\log X-\log Y)$. Let the weights and positive scalar
scales be fixed, and define
\begin{equation}\label{eq:qre-moments}
 A=\sum_iw_ic_i,
 \qquad B=\sum_iw_ic_i\log(c_i/d_i),
 \qquad D_{\mathrm{eff}}=A e^{-B/A}.
\end{equation}
Then
\begin{equation}\label{eq:qre-compiler}
 \sum_iw_iD(c_iX\|d_iY)
     =A D(X\|Y)+B\tr X
     =D(AX\|D_{\mathrm{eff}}Y).
\end{equation}
Indeed, scalar multiplication adds $(\log c)I$ to a matrix logarithm;
no commutativity of $X$ and $Y$ is required. The identity extends to the
positive-semidefinite closure with the usual condition
$\supp X\subseteq\supp Y$ for finite relative entropy. If that condition
fails, both sides are $+\infty$. If $X=0$, both sides vanish.

After acquisition of $A,B$, all value and derivative evaluations use
one relative-entropy term and no further access to the scenarios.
The scalar restriction, in which $-\log r$ has dilation rank two,
shows that two continuous real statistics are necessary for general
positive coefficient/scale families. If $A$ is fixed in advance, only
$B$ needs to be stored. For $c_i=1$ and $\sum_iw_i=1$,
$D_{\mathrm{eff}}=\prod_i d_i^{w_i}$.

The convention matters. For
$D_+(X\|Y)=D(X\|Y)+\tr Y-\tr X$ and
$C=\sum_iw_id_i$, the exact identity instead is
\[
 \sum_iw_iD_+(c_iX\|d_iY)
 =D_+(AX\|D_{\mathrm{eff}}Y)+(C-D_{\mathrm{eff}})\tr Y.
\]
The compiler preserves a weighted sum in one objective or epigraph;
it does not preserve separate scenario inequalities or separate outputs.
Nor does its source-free online evaluation apply when the scales change
with the optimization variables.

The relative-entropy epigraph barrier of parameter $2n+1$ for $n\times n$
matrices is established by \citet{FS2023}. Positive-map composition and
removal of redundant input log determinants are treated more generally
by \citet{HeSaundersonFawzi2026}. Equation~\eqref{eq:qre-compiler} is an elementary
scaling identity; its role here is to make an exact source summary
explicit. It is not a new barrier construction.

Acquiring the moments is a separate task. For a uniform average of
bounded coherently accessible entries, quantum amplitude estimation
uses $O(\min\{N,\eta^{-1}\})$ queries for constant-success additive
error $\eta$, compared with the randomized upper bound
$O(\min\{N,\eta^{-2}\})$ \cite{BHMT2002}. This statement presumes that
preparing the uniform index state and computing an entry are included
in the access model; arbitrary implicit weights are not free. On Boolean
entries, exact acquisition, or error strictly below $1/(2N)$ in the
mean, recovers the Hamming weight and requires $\Omega(N)$ quantum
queries \cite{NayakWu1999}. If $A$ is known exactly, an estimate
$\widehat B$ changes the epigraph boundary by precisely
$(\widehat B-B)\tr X$. Since $X\succeq0$, one-sided bounds on $B$
therefore give globally valid inner and outer epigraphs.

A small barrier parameter alone does not supply source-free evaluation.
For $z\in\{0,1\}^N$, let
\[
 \phi_z(x_1,x_2)=\sum_{i=1}^N
        D(x_1\|e^{-z_i}x_2)
       =N x_1\log(x_1/x_2)+|z|x_1.
\]
The epigraph on $x_1,x_2>0$ is a linear image of a single scalar
relative-entropy cone: use
$D(Nx_1\|Nx_2)=N x_1\log(x_1/x_2)$ and shift the epigraph coordinate
by $|z|x_1$. Thus
\[
 F_z(t,x_1,x_2)=-\log\bigl(t-\phi_z(x_1,x_2)\bigr)
                      -\log x_1-\log x_2
\]
is a three-parameter logarithmically homogeneous self-concordant barrier,
up to an additive constant \cite{FS2023}. Under the promise $|z|\in\{0,1\}$,
its value at the common interior point $(2,1,1)$ differs by exactly
$\log2$. An oracle returning that barrier value with error less than
$\log2/3$ solves unique-search decision and requires $\Omega(\sqrt N)$
quantum or $\Omega(N)$ randomized source queries. The lower bound is on
initial source access: it does not apply after $|z|$ has been learned and
stored.

\subsection{Exponential features versus exponential-cone factors}

Use the coordinate convention
\[
 K_{\exp}=\operatorname{cl}\{(a,b,c): b>0,\ c\geq b e^{a/b}\}.
\]
In particular, $(a,1,c)\in K_{\exp}$ means $c\geq e^a$. An exact
\emph{signed exponential representation} of a function on an open set is
an identity $f(z)=\sum_{j=1}^M\beta_j e^{\ip{\lambda_j}{z}}$ with real
nodes $\lambda_j$ and real coefficients $\beta_j$. Its rank is the
minimum number of nonzero terms, allowing arbitrary competing nodes.

\begin{proposition}[Unbounded feature rank with three cones]
\label{prop:binomial-three-exp}
For every integer $m\geq1$, the function
\[
 F_m(z)=\left(\frac{1+e^{z/m}}2\right)^m
\]
has signed exponential rank $m+1$ on every nonempty open interval, but
its full epigraph has an exact lift using three exponential cones and
one scalar linear inequality:
\begin{equation}\label{eq:binomial-exp-lift}
 (z/m-s,1,u),\ (-s,1,v),\ (ms,1,t)\in K_{\exp},
 \qquad u+v\leq2.
\end{equation}
All affine coefficients have $O(\log(m+1))$-bit rational encodings.
\end{proposition}
\begin{proof}
The first two cones and the inequality admit $u,v$ exactly when
$s\geq\log((1+e^{z/m})/2)$. The last cone then projects exactly to
$t\geq F_m(z)$, with equality attainable. On the other hand,
\[
 F_m(z)=2^{-m}\sum_{j=0}^m\binom mj e^{jz/m}.
\]
Distinct real exponentials are linearly independent on an open interval:
differentiating a finite zero combination at any interior point gives
a Vandermonde system. Comparing any competing representation with this
one therefore forces all $m+1$ displayed nodes to remain.
\end{proof}

Here $F_m$ is the moment-generating function of
$m^{-1}\sum_{j=1}^m B_j$ for independent Bernoulli$(1/2)$ variables,
so even restricting to probability measures supported on $[0,1]$ does
not restore a feature-rank bound in terms of the number of cones. The
log-sum-exp and outer-exponential operations used in
\eqref{eq:binomial-exp-lift} are standard modeling constructions
\cite[Sections~5.2.5--5.2.6]{MosekExpoCookbook}. The separation comes from
comparing this composition with the exact signed-feature output contract.

\begin{theorem}[Bounds for exact exponential lifts of a product-cosh graph]
\label{thm:cosh-exp-count}
Let $r\geq1$, $K>0$, and
$f_r(z)=\prod_{i=1}^r\cosh z_i$. Its signed exponential rank on every
nonempty open subset of $\R^r$ is $2^r$. Its epigraph has a lift with
$2r+1$ exponential-cone factors and $r$ scalar inequalities. Conversely,
any exact affine lift of its epigraph over $[-K,K]^r$ through
$K_{\exp}^R\times\R_+^P$, with arbitrary affine projections and free
auxiliary variables, satisfies
\begin{equation}\label{eq:cosh-exp-sandwich}
 R\geq r.
\end{equation}
Consequently the minimum exponential-cone factor count, with scalar
inequalities uncharged, lies between $r$ and $2r+1$. For any such
original ambient product, every self-concordant barrier has parameter
at least $3R+P\geq3r+P$, and the ambient optimum is $3R+P$.
\end{theorem}
\begin{proof}
The expansion
$f_r(z)=2^{-r}\sum_{\epsilon\in\{-1,1\}^r}e^{\ip{\epsilon}{z}}$
has $2^r$ distinct nodes. To prove independence against any competing
finite node set, choose a direction not perpendicular to any pairwise
node difference, and restrict the identity to a short line segment
inside the given open set. The one-dimensional Vandermonde argument
then applies.

For the upper bound impose, for each $i$,
\[
 (z_i-s_i,1,u_i),\ (-z_i-s_i,1,v_i)\in K_{\exp},
 \quad u_i+v_i\leq2,
 \qquad \left(\sum_i s_i,1,t\right)\in K_{\exp}.
\]
Eliminating $u_i,v_i$ gives $s_i\geq\log\cosh z_i$, and eliminating
$s$ gives exactly $t\geq f_r(z)$. Restricting to the cube adds only
scalar inequalities.

For the lower bound choose $T>1$ and intersect the given epigraph with
$t\leq T$. Its image is the compact convex body
\[
 C=\{(z,t):z\in[-K,K]^r,\ f_r(z)\leq t\leq T\}
       \subset\R^{r+1}.
\]
It has nonempty interior. In a neighborhood of $(0,1)$ its lower
boundary is the graph of $f_r$, and
\begin{equation}\label{eq:cosh-hessian}
 \nabla^2 f_r(z)=f_r(z)\left[
    \operatorname{diag}(\operatorname{sech}^2 z_i)
                  +(\tanh z)(\tanh z)^T\right]\succ0.
\end{equation}
Its second fundamental form is therefore positive definite in all
$r$ tangent directions. The intersection adds one nonnegative-ray
factor and does not add an exponential factor. After translation to
put an interior point at the origin, Lemma~\ref{lem:reduction}
eliminates free variables and passes to the minimal face. These sets
are definable in the real exponential structure; the generic selection
argument and Theorem~\ref{thm:curvature} apply. Each retained face of an
exponential cone has dimension at most three and thus contributes at
most one unit to $\sum_i\max\{\dim K_i-2,0\}$; every ray contributes
zero. The $r$-dimensional positive-curvature patch gives $r\leq R$.
Finally, Theorem~\ref{thm:entropy-product} gives the exact ambient
parameter $3R+P$ for the \emph{original} product. The artificial
truncation ray is used only in the curvature argument and is not charged
to that original product.
\end{proof}

The same proof gives $R\geq r$ for any definable convex epigraph with
an $r$-dimensional $C^2$ graph patch having positive definite Hessian,
provided that a compact polyhedral truncation retains that patch.
It is important to keep the epigraph coordinate in this argument:
a level-set section of $f_r$ alone has only $r-1$ curved directions.
Theorem~\ref{thm:cosh-exp-count} makes no optimality claim for the
constant $2$ in its constructive upper bound. Its lower bound follows
from curvature, independently of exponential-feature rank.

\begin{proposition}[Polyhedral obstruction to an approximate transfer]
\label{prop:exp-polyhedral-approximation}
Let $D\subset\R^r$ be a compact polytope and let $f>0$ be convex and
$C^2$ on a neighborhood of $D$. For every $0<\varepsilon<1$, a finite
piecewise-affine convex function $p$ satisfies
\[
 (1-\varepsilon)f\leq p\leq f\quad\hbox{on }D.
\]
In particular,
\[
 \operatorname{epi}\bigl(p/(1-\varepsilon)\bigr)
 \subseteq\operatorname{epi} f\subseteq\operatorname{epi}p
 \quad\hbox{over }D,
\]
and the two approximating epigraphs use only scalar linear inequalities.
\end{proposition}
\begin{proof}
Let $m_f=\min_D f>0$ and $L_f=\sup_D\norm{\nabla^2 f}_{\mathrm{op}}$.
If $L_f=0$, one tangent is exact on $D$. Otherwise choose a finite
$\delta$-net $\mathcal N\subset D$, with
$\delta^2\leq2\varepsilon m_f/L_f$, and set
\[
 p(z)=\max_{x\in\mathcal N}
          \{f(x)+\ip{\nabla f(x)}{z-x}\}.
\]
Convexity gives $p\leq f$. Taylor's theorem on the segment between $z$
and a net point within $\delta$ gives
$f(z)-p(z)\leq L_f\delta^2/2\leq\varepsilon f(z)$.
The epigraph of a finite maximum of affine functions is their halfspace
intersection, giving the claim.
\end{proof}

Thus uniform relative approximation on a fixed cube can use zero
exponential cones when arbitrarily many scalar inequalities are free.
Charging those inequalities, or imposing the conditioned contact
hypotheses of the preceding approximation results, changes the question.
Neither exact signed-feature rank nor the exact curvature count alone
is an approximation lower bound. Together, the three compiler models
identify which representation or access restriction supplies each
conclusion; none yields a lower bound on all optimization algorithms.
